\subsection{Metodología de Desarrollo: Modelo en Espiral}

El Modelo en Espiral, propuesto por Barry W. Boehm en 1986, es un modelo de proceso
de software de naturaleza iterativa e incremental que combina elementos del modelo en
cascada con el prototipado, incorporando de manera explícita la gestión de riesgos en
cada ciclo de desarrollo. A diferencia de los modelos lineales, la espiral no asume que
los requisitos son completamente conocidos desde el inicio, sino que permite refinarlos
y ampliarlos en cada vuelta del ciclo. \cite{boehm1988spiral}

Cada iteración del modelo sigue cuatro cuadrantes fundamentales:

\begin{enumerate}
    \item \textbf{Determinación de objetivos, alternativas y restricciones.} Se definen
    las metas del ciclo, se identifican las alternativas de solución y se establecen las
    restricciones técnicas, económicas y de tiempo.

    \item \textbf{Evaluación de alternativas e identificación y resolución de riesgos.}
    Se analizan los posibles riesgos asociados a cada alternativa y se diseñan estrategias
    de mitigación antes de comprometer recursos.

    \item \textbf{Desarrollo y verificación del producto.} Se implementa el incremento
    correspondiente al ciclo y se valida mediante pruebas y revisiones.

    \item \textbf{Planificación de la siguiente fase.} Con base en los resultados
    obtenidos, se planifica la siguiente iteración y se ajustan objetivos y alcance de la próxima iteración.
\end{enumerate}

\subsubsection{Justificación de la elección}

La metodología en Espiral es la más adecuada para este proyecto debido a la naturaleza
de sus componentes de Inteligencia Artificial y Visión por Computadora, los cuales
requieren iteraciones de entrenamiento, validación y ajuste que una metodología lineal
no permite gestionar eficazmente. En cada vuelta de la espiral se construye un
incremento del sistema, se prueba y se reduce el riesgo acumulado.

Entre las razones específicas que justifican esta elección destacan:

\begin{itemize}
    \item \textbf{Alta incertidumbre técnica.} La integración de OCR especializado en
    escritura infantil en español y de módulos de NLP implica decisiones tecnológicas —
    selección de modelos, datasets, estrategias de \textit{fine-tuning} — cuyo resultado
    no puede garantizarse desde el inicio del proyecto.

    \item \textbf{Gestión explícita de riesgos.} Dado que el proyecto presenta riesgos en múltiples dimensiones: técnicos (precisión del OCR
    ante trazos irregulares o inmaduros), operativos (adopción por parte de docentes),
    legales (protección de datos de menores) y de disponibilidad (dependencia de APIs
    externas). El modelo en Espiral sitúa la evaluación de riesgos como actividad
    central de cada ciclo.

    \item \textbf{Requisitos evolutivos.} El sistema atiende a un público con necesidades
    pedagógicas específicas, por lo que los requisitos funcionales y de usabilidad debían
    validarse progresivamente conforme avanzaba el levantamiento de información con
    representantes del dominio educativo.

    \item \textbf{Desarrollo incremental por módulos.} La arquitectura contempla
    componentes claramente diferenciados — autenticación, captura de escritura, OCR,
    NLP y retroalimentación — que pueden desarrollarse y validarse de manera independiente
    en iteraciones sucesivas.
\end{itemize}

\subsubsection{Estructura de las iteraciones}

El desarrollo se estructuró en cuatro ciclos evolutivos distribuidos en las dos etapas
del Trabajo Terminal. Las iteraciones 1 y 2 (correspondientes a TT1) se enfocaron en
la construcción de la infraestructura web, la recolección del dataset y el
preprocesamiento de imágenes, mitigando tempranamente los riesgos asociados a la
calidad de la entrada. Las iteraciones 3 y 4 (correspondientes a TT2) abordarán la
implementación de los modelos de Inteligencia Artificial (OCR y NLP) y la validación
final con usuarios, asegurando que los componentes críticos se desarrollen sobre una
base de datos validada y estable.

\paragraph{Etapa 1 — Trabajo Terminal I: Fundamentos y Procesamiento Visual}

El objetivo de esta etapa fue establecer la infraestructura tecnológica, recolectar el
corpus de datos necesario y validar los algoritmos de preprocesamiento de imágenes.

\textbf{Iteración 1: Infraestructura Web y Módulo de Captura}

\textit{Objetivo:} Desarrollar la arquitectura base de la aplicación web (Frontend y
Backend) y el módulo de digitalización de escritura.

\textit{Actividades principales:}
\begin{itemize}
    \item Diseño y creación de la base de datos relacional.
    \item Implementación de mecanismos de autenticación y gestión de usuarios.
    \item Desarrollo de la interfaz de carga y captura de imágenes.
\end{itemize}

\textit{Gestión de riesgos:}
\begin{itemize}
    \item \textbf{Riesgo:} Variabilidad en la calidad de las imágenes capturadas
    (iluminación, resolución) que impida su procesamiento futuro.
    \item \textbf{Mitigación:} Implementación de validadores de calidad en el cliente
    y normalización básica de imágenes al momento de la subida o permitir subir por el lienzo digital.
\end{itemize}

\textbf{Iteración 2: Conformación del Conjunto de Datos y Preprocesamiento}

\textit{Objetivo:} Seleccionar y adaptar datasets existentes de escritura manuscrita
con letra de molde e incorporar muestras adicionales, incluyendo caracteres con tilde,
para completar la variabilidad necesaria en el entrenamiento y validación de los modelos
de reconocimiento.

\textit{Actividades principales:}
\begin{itemize}
    \item Recolección y etiquetado de muestras de escritura manuscrita de tercer y
    cuarto grado.
    \item Implementación de algoritmos de binarización, eliminación de ruido y
    corrección de perspectiva.
    \item Desarrollo de algoritmos de segmentación para identificar renglones y palabras
    individuales.
\end{itemize}

\textit{Gestión de riesgos:}
\begin{itemize}
    \item \textbf{Riesgo:} Insuficiencia de datos públicos en español para el
    entrenamiento de modelos.
    \item \textbf{Mitigación:} Creación de un dataset propio validado para integrarlo
    con un conjunto existente, asegurando la diversidad caligráfica necesaria para el
    siguiente ciclo.
\end{itemize}

\paragraph{Etapa 2 — Trabajo Terminal II: Inteligencia Artificial e Integración}

En esta etapa se integran los modelos de aprendizaje profundo y procesamiento de
lenguaje natural, utilizando los datos procesados en la etapa anterior.

\textbf{Iteración 3: Implementación de Modelos OCR y Métricas Caligráficas}

\textit{Objetivo:} Integrar el motor de Reconocimiento Óptico de Caracteres (OCR) y
los algoritmos de evaluación de calidad del trazo, así como la generación del módulo
de retroalimentación mediante ejercicios correctivos.

\textit{Actividades principales:}
\begin{itemize}
    \item Entrenamiento y \textit{fine-tuning} del modelo OCR (TrOCR/Tesseract)
    utilizando el corpus generado en la Iteración 2.
    \item Programación de métricas geométricas para evaluar inclinación, espaciado y
    tamaño de letra.
    \item Validación técnica de la precisión de la transcripción.
    \item Desarrollo del módulo de retroalimentación con ejercicios correctivos.
\end{itemize}

\textit{Gestión de riesgos:}
\begin{itemize}
    \item \textbf{Riesgo:} Baja precisión del modelo ante trazos irregulares o
    inmaduros.
    \item \textbf{Mitigación:} Reentrenamiento iterativo y ajuste de hiperparámetros
    específicos para la morfología de la letra infantil.
\end{itemize}

\textbf{Iteración 4: Análisis Lingüístico (NLP), Integración y Despliegue}

\textit{Objetivo:} Implementar la detección semántica de errores ortográficos y la
generación de reportes pedagógicos, culminando con el despliegue del sistema completo.

\textit{Actividades principales:}
\begin{itemize}
    \item Desarrollo del módulo de NLP para la corrección ortográfica contextual y
    gramatical.
    \item Generación automática de retroalimentación para el alumno y reportes de
    progreso para el docente.
    \item Pruebas de aceptación con usuarios finales y despliegue del sistema.
\end{itemize}

\textit{Gestión de riesgos:}
\begin{itemize}
    \item \textbf{Riesgo:} Generación de falsos positivos en la corrección ortográfica.
    \item \textbf{Mitigación:} Calibración de umbrales de confianza y validación
    cruzada con evaluaciones docentes manuales.
\end{itemize}

La Tabla~\ref{tab:espiral} resume la correspondencia entre cada iteración, su etapa
del Trabajo Terminal y los entregables principales asociados.

\begin{table}[H]
\centering
\caption{Correspondencia entre iteraciones del modelo en espiral y entregables del proyecto.}
\label{tab:espiral}
\begin{tabular}{|p{1.5cm}|p{2.5cm}|p{3.5cm}|p{5cm}|}
\hline
\textbf{Iteración} & \textbf{Etapa} & \textbf{Enfoque} & \textbf{Entregables principales} \\
\hline
1 & TT1 & Infraestructura web y captura & Base de datos, autenticación, módulo de captura de imágenes \\
\hline
2 & TT1 & Dataset y preprocesamiento & Corpus etiquetado, algoritmos de binarización y segmentación \\
\hline
3 & TT2 & OCR y métricas caligráficas & Modelo OCR ajustado, métricas geométricas, módulo de retroalimentación \\
\hline
4 & TT2 & NLP, integración y despliegue & Módulo NLP, reportes de progreso, sistema integrado desplegado \\
\hline
\end{tabular}
\end{table}